% ===================================================================
% CHAPTER 6: SYSTEM ARCHITECTURE
% ===================================================================
\chapter{Архитектурна логика на системот}
\label{chap:system}

Ова поглавје го опишува системот што ги произведе мерењата од \hyperref[chap:results]{Дел~\ref*{chap:results}}. Описот е
намерно конкретен --- со вистинските имиња на услуги, променливи на околината и датотеки --- затоа што репродуцибилноста на физичко мерење зависи од детали што во апстрактен дијаграм
се губат.

\section{Податочен извор и негова употреба}

Системот работи врз податочното множество TON\_IoT \cite{moustafa2021toniot}, објавено од
UNSW, конкретно врз неговите 23 CSV-датотеки со означен мрежен сообраќај (≈ 3,3 GB). Тие
датотеки играат две улоги истовремено, што е важно за разбирањето на дизајнот:

\begin{enumerate}
  \item Тие се \textbf{изворот на федеративната хетерогеност}. Поделбата меѓу јазлите се врши на ниво на \emph{список од датотеки}, а не на ниво на редови, па секој јазол добива дисјунктно множество целосни снимки. Резултатот е природна non-IID распределба по извор на сообраќај, без вештачки внесена пристрасност.
  \item Тие се \textbf{изворот на класната нерамнотежа} што ја определува целата расправа за корисноста. Уделот на напади во множеството се наследува од самото множество, не се израмнува, и изнесува околу 96,5\,\%. Тоа е реалистично за означено множество наменето за откривање упади, и токму затоа изборот на показатели за оценување е одлучувачки (\hyperref[chap:theory]{Дел~\ref*{chap:theory}}, §3.5).
\end{enumerate}

\section{Детален системски дизајн и податочен проток}

\subsection{Рамка и режим на работа}

Системот е изграден врз Flower 1.29.0 \cite{beutel2020flower} во \textbf{режим на распоредување}
(\emph{deployment}), а не во режим на симулација. Разликата е суштинска за овој труд: во режим на
распоредување секој учесник е вистински процес во сопствен контејнер, што комуницира преку
мрежа, па измереното време и енергија ја вклучуваат и реалната режија на комуникацијата и
на распоредувањето.

Компонентите се:

\begin{table}[htbp]
\centering
\caption{Компоненти на топологијата во режим на распоредување и нивната улога.}
\label{tab:topology-components}
\begin{tabular}{lp{0.68\linewidth}}
\toprule
Компонента & Улога \\
\midrule
\texttt{superlink} & Централен координатор и посредник за пораки. Изложува два интерфејса: Fleet API (порта 9092), на кој се приклучуваат јазлите, и Exec API (порта 9093), преку кој се предава работата. \\
\texttt{supernode-1..4} & По еден на секој јазол. Секој се стартува со \texttt{--node-config "partition-id=N num-partitions=4"} --- така јазолот дознава која партиција е негова. \\
\texttt{superexec-clientapp-1..4} & Приклучоци за извршување, врзани за соодветниот \texttt{supernode}, кои го извршуваат \texttt{edge\_nodes/client\_app.py}. \\
\texttt{superexec-serverapp} & Приклучок што го извршува \texttt{server/server\_app.py}; стратегијата, бројот на рунди и минималниот број клиенти се читаат од променливи на околината. \\
\texttt{preprocessor} & Еднократна задача што ги гради кешовите по партиција; сите клиентски контејнери зависат од нејзиното успешно завршување. \\
\texttt{runner} & Контејнер за еднократна употреба што ја предава работата со \texttt{flwr run .} и ги пренесува дневниците. \\
\texttt{toxiproxy} & Посредник меѓу јазлите и \texttt{superlink} за внесување мрежно доцнење. \\
\bottomrule
\end{tabular}
\end{table}

Врзувањето на кодот за рамката е во \texttt{pyproject.toml}, во \texttt{[tool.flwr.app.components]}:
\texttt{serverapp = "server.server\_app:app"} и \texttt{clientapp = "edge\_nodes.client\_app:app"}.

\subsection{Хибридна топологија}

Мерниот експеримент користи \textbf{хибридна} поставка распоредена преку две композициски
датотеки:

\begin{itemize}
  \item \textbf{\texttt{docker-compose.host.yml}} се извршува на домаќинот (лаптоп) и подига \texttt{superlink}, \texttt{superexec-serverapp} и симулираните јазли 1–3. Fleet API (порта 9092) е изложен кон локалната мрежа за да може физичкиот јазол да се приклучи. Трите симулирани клиенти добиваат \textbf{намерно нееднакви} ограничувања --- 3,0 / 2,0 / 1,5 процесорски јадра и 3 GB / 2 GB / 1500 MB меморија --- со што федерацијата е хетерогена, како во реална поставка.
  \item \textbf{\texttt{docker-compose.edge.yml}} се извршува на Raspberry Pi 5 и подига \texttt{supernode-4} (партиција 3 од 4) и неговиот клиентски приклучок. Двата контејнери работат со \texttt{network\_mode: host} и се поврзуваат кон \texttt{HOST\_IP:9092}.
\end{itemize}

Две одлуки во конфигурацијата на јазолот се директно во функција на мерењето:

\begin{itemize}
  \item Директориумот со контролни точки е монтиран \textbf{на SD-картичката} (\texttt{CHECKPOINT\_HOST\_DIR}), намерно, бидејќи влезно-излезниот трошок на контролните точки е една од мерените величини.
  \item Телеметријата се запишува \textbf{исклучиво во \texttt{/dev/shm}} (RAM), за метриката за запишување на SD-картичката да го одразува само работното оптоварување, а не сопственото логирање.
\end{itemize}

Понашањето на клиентот се менува преку две променливи: \texttt{CLIENT\_MODE} (\texttt{standard} за
пристапот А, \texttt{sisa} за пристапот Б) и \texttt{POISON\_MODE} (\texttt{flip} во Фаза 1, \texttt{drop} во Фаза 4,
\texttt{off} за чисто извршување).

\subsection{Внесување мрежни услови}

\texttt{toxiproxy} е поставен меѓу јазлите и \texttt{superlink}. Самиот посредник е декларативно опишан во
\texttt{toxiproxy.json}, но токсинот (доцнењето) се додава со REST-повик при стартувањето, бидејќи
статичкиот формат на toxiproxy не може да изрази токсини. Во симулацијата од претходната
фаза доцнењето е постојано вклучено и изнесува 5 ms ± 3 ms; во магистерскиот експеримент тоа
е \textbf{секундарен фактор} зад профилот \texttt{wan} (40 ms ± 20 ms), додека примарните мерења се
изведени врз чиста локална мрежа.

\subsection{Агрегација на страната на серверот}

\texttt{server/server\_app.py} избира стратегија според \texttt{FL\_STRATEGY}: \texttt{fedavg} (сопствената
\texttt{ResilientFedAvg}), \texttt{fedprox}, \texttt{krum} или \texttt{trimmedmean}. Во секоја рунда серверот ги
испраќа хиперпараметрите (\texttt{local\_epochs}, \texttt{lr}) кон клиентите, што значи дека промена на
хиперпараметар не бара повторно градење на клиентските слики.

Едно својство на \texttt{ResilientFedAvg} заслужува посебно спомнување: ако во дадена рунда
одговорат помалку клиенти од пропишаниот минимум, таа ги \textbf{задржува тежините од претходната
рунда} наместо да агрегира нецелосно ажурирање. Без тоа, привремен прекин на врската со
физичкиот јазол --- во хибридна поставка преку вистинска мрежа, сосема очекуван настан --- би
можел да ја разниша обуката и да внесе варијабилност што нема врска со предметот на
мерењето.

\section{Процесирање и трансформација на податоци}

Податочниот проток има три чекори, изведени со три различни модули.

\subsection{Чекор 1 --- градење кешови по партиција (\texttt{edge\_nodes/preprocess.py})}

Модулот се извршува еднаш врз сите 23 датотеки:

\begin{enumerate}
  \item Се отфрлаат метаподатоците и колоните со етикети (\texttt{ts}, \texttt{src\_ip}, \texttt{dst\_ip}, \texttt{src\_mac}, \texttt{dst\_mac}, \texttt{type}).
  \item Се пресметува \textbf{пресекот на нумеричките колони заеднички за сите датотеки}, што дава доследен вектор на обележја со димензија \textbf{15}; списокот се запишува во \texttt{columns.json}.
  \item Списокот со датотеки (не редовите) се дели на \texttt{NUM\_PARTITIONS} делови.
  \item Секоја партиција се стандардизира \textbf{независно} (нулта средна вредност, единечна дисперзија) и се запишува како \texttt{data/.cache/partition\_\{i\}\_of\_\{N\}.npz}.
\end{enumerate}

Постапката е идемпотентна: ако кешовите и \texttt{columns.json} веќе постојат, се прескокнува.
Магистерскиот експеримент бара кешови за четири партиции и затоа се гради со
\texttt{NUM\_PARTITIONS=4}.

\subsection{Чекор 2 --- подмножество, тест-множество и труење (\texttt{experiments/prepare\_edge\_data.py})}

Од секоја партиција се зема \textbf{стратификувано} подмножество од 1\,000 000 реда за обука,
еднакво за сите четири јазли --- намерно, за придонесот на секој клиент под FedAvg да биде
еднакво тежински вреднуван. Одделно се издвојуваат по 25\,000 реда по партиција, дисјунктни
од множеството за обука, кои заедно го формираат \textbf{глобалното тест-множество од 100\,000
реда}. Тоа множество се оценува на страната на домаќинот и е основата за секое тврдење за
корисноста во \hyperref[chap:results]{Дел~\ref*{chap:results}}.

Труењето се внесува само за партицијата 3 (физичкиот јазол). Пресметката на индексите користи
\textbf{одделен случаен тек} (\texttt{seed + 1}) од оној што го распоредува множеството по шардови, а
изборот се ограничува на нападите во сегментите со индекс ≥ 3 од шардот 1.

Најважниот детаљ во оваа конструкција: \textbf{етикетите на дискот остануваат чисти}. Се чува
само низата индекси \texttt{poison\_idx}, а превртувањето се изведува во меморија при вчитувањето,
кога \texttt{POISON\_MODE=flip}. Тоа значи дека двата крака работат врз \textbf{битово идентично}
податочно множество, и дека обновувањето може едноставно да ги испушти тие редови и да ги
користи вистинските етикети. Истата низа индекси ја игра улогата на \textbf{оракул-маска} во Фаза 2.

\subsection{Чекор 3 --- вчитување на јазолот (\texttt{edge\_nodes/data\_loader.py})}

Вчитувањето применува \textbf{позициска} поделба 80/20 на обука и валидација (редовите се веќе
пермутирани во чекорот 2, па повторно мешање не е потребно). Трите режима се \texttt{off},
\texttt{flip} (превртување во меморија) и \texttt{drop} (\texttt{np.setdiff1d} врз индексите за обука).
Декомпримираните низи се задржуваат во RAM меѓу повиците, за да не се чита повторно од
SD-картичката во секоја рунда --- уште една мерка за метриката за влезно-излезни операции да
го одразува само работното оптоварување.

\subsection{Моделот}

\texttt{edge\_nodes/model.py} дефинира повеќеслоен перцептрон со димензии
$15 \to 128 \to 64 \to 32 \to 1$. Секој скриен слој се состои од линеарна трансформација,
активација ReLU, пакетна нормализација и испуштање со $p = 0{,}3$; излезниот неврон користи
сигмоидна активација. Обуката користи \texttt{BCELoss} и оптимизаторот Adam со $lr = 0{,}001$ и
големина на пакет 512. Моделот има \textbf{12\,865} параметри што се учат --- податок што подоцна се
покажува како одлучувачки за исходот на хипотезата H4.

\subsection{Слојот SISA}

Три модули го сочинуваат:

\begin{itemize}
  \item \textbf{\texttt{sisa\_partition.py}} ја гради распределбата: пермутација определена со почетната вредност, поделена на $S = 5$ шардови, секој на $R = 5$ сегменти. Функцијата е чиста --- зависи само од бројот примери и од почетната вредност --- па истата распределба се реконструира идентично и при подготовката на податоците, и при обуката, и при обновувањето.
  \item \textbf{\texttt{sisa\_client.py}} одржува по еден составен модел на шард. По секој сегмент запишува контролна точка што ја содржи и состојбата на моделот и \textbf{состојбата на оптимизаторот} (за да може обновувањето да продолжи точно од средината на рундата); секоја контролна точка има големина од 174\,688 бајти. Две особености се документирани отстапувања од класичниот SISA и имаат директни последици: составните модели \textbf{никогаш не ги преземаат глобалните тежини} (што е условот гаранцијата при враќање наназад да важи), а ажурирањето што се испраќа кон FedAvg е \textbf{параметарска средна вредност} на петте составни модели, наместо ансамбл од предвидувања.
  \item \textbf{\texttt{sisa\_recover.py}} го изведува обновувањето: за секој шард го наоѓа \textbf{првиот} компромитиран сегмент, се враќа на контролната точка од претходниот сегмент од рундата 1, и ги обработува повторно само сегментите од таа точка наваму. Во оваа поставка тоа е секогаш \texttt{\{шард 1: сегмент 3\}} и \textbf{47 сегменти} за повторна обработка, наспроти 250 за целосна обука. Обновениот модел е средна вредност на петте составни модели.
\end{itemize}

\subsection{Наивното обновување}

\texttt{edge\_nodes/naive\_retrain.py} е намерно едноставен: ги отфрла компромитираните редови
(остануваат ≈ 736\,940 од 800\,000), иницијализира \textbf{целосно нов} модел и го обучува 10
епохи со нов оптимизатор во секоја епоха. Буџетот од 10 епохи не е произволен --- тој е
еднаков на локалниот буџет од Фаза 1 (10 рунди × 1 локална епоха), па двата крака добиваат
споредлива пресметковна дозвола.

\section{Развојна околина и технологии}

\subsection{Верзии}

\begin{table}[htbp]
\centering
\caption{Верзии на софтверот употребен во експериментот.}
\label{tab:versions}
\begin{tabular}{ll}
\toprule
Компонента & Верзија \\
\midrule
Flower & 1.29.0 (сите слики: \texttt{superlink}, \texttt{supernode}, \texttt{superexec}, \texttt{base}) \\
Python & 3.13 \\
PyTorch & 2.13.0 \\
NumPy & 2.4.4 \\
pandas & 3.0.2 \\
scikit-learn & 1.8.0 \\
Основна слика на контејнерите & Ubuntu 24.04 \\
Оперативен систем на јазолот & Raspberry Pi OS Lite, Debian 13 „trixie“, 64-битен \\
\bottomrule
\end{tabular}
\end{table}

\subsection{Подготовка на јазолот (\texttt{ansible/setup\_node.yaml})}

Поставувањето е автоматизирано, и најголемиот дел од него постои заради \textbf{чистотата на
мерењето}, не заради функционалноста:

\begin{itemize}
  \item Се запираат периодичните задачи за ажурирање (\texttt{apt-daily.timer}), кои инаку создаваат случајни скокови во оптоварувањето на процесорот и во влезно-излезните операции.
  \item Виртуелната меморија се оневозможува трајно. На Debian 13 тоа бара запис \texttt{Mechanism=none} во \texttt{/etc/rpi/swap.conf}, бидејќи новиот механизам \texttt{rpi-swap} го заменил \texttt{dphys-swapfile} и без ова виртуелната меморија тивко се враќа по рестарт --- и ја загадува токму метриката што се мери.
  \item Регулаторот на фреквенцијата се поставува на \texttt{performance}, за да не се меша автоматското скалирање со термичкото придушување што се набљудува.
  \item Дневниците на Docker се ограничуваат по големина, повторно за да не создаваат позадинско запишување на SD-картичката.
  \item Се поставува \texttt{monitor.sh}, телеметриската скрипта.
\end{itemize}

\subsection{Телеметрија и усогласување на часовниците}

Мерењето се потпира на три тека што мора да се поврзат:

\begin{table}[htbp]
\centering
\caption{Телеметриски текови и нивниот часовник и фреквенција на земање примероци.}
\label{tab:telemetry-streams}
\begin{tabular}{llll}
\toprule
Тек & Извор & Часовник & Фреквенција \\
\midrule
Хардверска телеметрија & \texttt{monitor.sh} на јазолот, пишува во \texttt{/dev/shm} & часовникот на јазолот & 1 Hz \\
Потрошувачка & \texttt{experiments/fnirsi\_logger.py} на домаќинот, преку USB од мерачот FNB58 & часовникот на домаќинот & $\approx$100 примероци/s \\
Фазни маркери & \texttt{phases.log} на домаќинот & часовникот на домаќинот & по настан \\
\bottomrule
\end{tabular}
\end{table}

\texttt{monitor.sh} во секоја секунда ги запишува температурата, работната фреквенција, регистарот
за придушување, искористената меморија, времето на чекање за влез/излез и пропусноста на
SD-картичката. Два детаља во неа се резултат на конкретни грешки избегнати во пилотот:
времето на чекање за влез/излез се пресметува како \textbf{разлика меѓу два последователни
отчита} од \texttt{/proc/stat} (инаку би се известувал просек од подигањето на системот наваму), а
пропусноста на SD-картичката се чита директно од бројачот на блок-уредот
(\texttt{/sys/block/mmcblk0/stat}), не од сметководството на датотечниот систем.

Усогласувањето на фазите е решено на два начина истовремено. Дневникот на потрошувачката и
фазните маркери \textbf{го делат часовникот на домаќинот}, така што интеграцијата на енергијата
по фаза не бара корекција за разлика меѓу часовници. Телеметријата на јазолот, пак, носи
\textbf{сопствена колона \texttt{Marker}} --- при секој премин меѓу фази домаќинот запишува ознака и во
\texttt{phases.log} и во \texttt{/dev/shm/run\_marker} на јазолот, а телеметриската јамка ја чита таа
датотека во секој отчит. Припаѓањето на отчитите кон фазите на страната на јазолот затоа се
утврдува по ознака, а не со аритметика врз временски печати.

\subsection{Проверка на исправноста}

\texttt{tests/smoke\_test.py} е единствениот автоматски тест во складиштето и служи за една
конкретна цел --- да ја поткрепи тврдењето за егзактност. Тој проверува дека распределбата по
шардови и сегменти и изборот на компромитирани индекси се детерминистички; дека режимот
\texttt{flip} ги менува етикетите само во меморија, а не на дискот; дека по обуката постои
контролна точка за секој сегмент; и --- најважно --- дека \textbf{двојно извршување на обновувањето
дава битово идентичен модел}, споредено преку SHA-256 отпечаток на сортираните тензори од
состојбата.

Опсегот на таа проверка треба да се наведе точно: таа докажува дека обновувањето е
репродуцибилно и дека точката на враќање наназад е избрана правилно. Таа \textbf{не} е
битова споредба со независно преобучен модел. Егзактноста на заборавањето почива на
конструкцијата --- отстранетите редови никогаш не влегуваат во повторно обработените
составни модели --- а тестот потврдува дека таа конструкција се однесува детерминистички.

